% Options for packages loaded elsewhere
% Options for packages loaded elsewhere
\PassOptionsToPackage{unicode}{hyperref}
\PassOptionsToPackage{hyphens}{url}
\PassOptionsToPackage{dvipsnames,svgnames,x11names}{xcolor}
%
\documentclass[
]{article}
\usepackage{xcolor}
\usepackage{amsmath,amssymb}
\setcounter{secnumdepth}{5}
\usepackage{iftex}
\ifPDFTeX
  \usepackage[T1]{fontenc}
  \usepackage[utf8]{inputenc}
  \usepackage{textcomp} % provide euro and other symbols
\else % if luatex or xetex
  \usepackage{unicode-math} % this also loads fontspec
  \defaultfontfeatures{Scale=MatchLowercase}
  \defaultfontfeatures[\rmfamily]{Ligatures=TeX,Scale=1}
\fi
\usepackage{lmodern}
\ifPDFTeX\else
  % xetex/luatex font selection
\fi
% Use upquote if available, for straight quotes in verbatim environments
\IfFileExists{upquote.sty}{\usepackage{upquote}}{}
\IfFileExists{microtype.sty}{% use microtype if available
  \usepackage[]{microtype}
  \UseMicrotypeSet[protrusion]{basicmath} % disable protrusion for tt fonts
}{}
\makeatletter
\@ifundefined{KOMAClassName}{% if non-KOMA class
  \IfFileExists{parskip.sty}{%
    \usepackage{parskip}
  }{% else
    \setlength{\parindent}{0pt}
    \setlength{\parskip}{6pt plus 2pt minus 1pt}}
}{% if KOMA class
  \KOMAoptions{parskip=half}}
\makeatother
% Make \paragraph and \subparagraph free-standing
\makeatletter
\ifx\paragraph\undefined\else
  \let\oldparagraph\paragraph
  \renewcommand{\paragraph}{
    \@ifstar
      \xxxParagraphStar
      \xxxParagraphNoStar
  }
  \newcommand{\xxxParagraphStar}[1]{\oldparagraph*{#1}\mbox{}}
  \newcommand{\xxxParagraphNoStar}[1]{\oldparagraph{#1}\mbox{}}
\fi
\ifx\subparagraph\undefined\else
  \let\oldsubparagraph\subparagraph
  \renewcommand{\subparagraph}{
    \@ifstar
      \xxxSubParagraphStar
      \xxxSubParagraphNoStar
  }
  \newcommand{\xxxSubParagraphStar}[1]{\oldsubparagraph*{#1}\mbox{}}
  \newcommand{\xxxSubParagraphNoStar}[1]{\oldsubparagraph{#1}\mbox{}}
\fi
\makeatother


\usepackage{longtable,booktabs,array}
\usepackage{calc} % for calculating minipage widths
% Correct order of tables after \paragraph or \subparagraph
\usepackage{etoolbox}
\makeatletter
\patchcmd\longtable{\par}{\if@noskipsec\mbox{}\fi\par}{}{}
\makeatother
% Allow footnotes in longtable head/foot
\IfFileExists{footnotehyper.sty}{\usepackage{footnotehyper}}{\usepackage{footnote}}
\makesavenoteenv{longtable}
\usepackage{graphicx}
\makeatletter
\newsavebox\pandoc@box
\newcommand*\pandocbounded[1]{% scales image to fit in text height/width
  \sbox\pandoc@box{#1}%
  \Gscale@div\@tempa{\textheight}{\dimexpr\ht\pandoc@box+\dp\pandoc@box\relax}%
  \Gscale@div\@tempb{\linewidth}{\wd\pandoc@box}%
  \ifdim\@tempb\p@<\@tempa\p@\let\@tempa\@tempb\fi% select the smaller of both
  \ifdim\@tempa\p@<\p@\scalebox{\@tempa}{\usebox\pandoc@box}%
  \else\usebox{\pandoc@box}%
  \fi%
}
% Set default figure placement to htbp
\def\fps@figure{htbp}
\makeatother





\setlength{\emergencystretch}{3em} % prevent overfull lines

\providecommand{\tightlist}{%
  \setlength{\itemsep}{0pt}\setlength{\parskip}{0pt}}



 


\usepackage{amsthm}
\newtheorem{task}{Task}
\newtheorem*{use case}{Use Case}
\newtheorem{problem}{Problem}[section]
\makeatletter
\@ifpackageloaded{caption}{}{\usepackage{caption}}
\AtBeginDocument{%
\ifdefined\contentsname
  \renewcommand*\contentsname{Table of contents}
\else
  \newcommand\contentsname{Table of contents}
\fi
\ifdefined\listfigurename
  \renewcommand*\listfigurename{List of Figures}
\else
  \newcommand\listfigurename{List of Figures}
\fi
\ifdefined\listtablename
  \renewcommand*\listtablename{List of Tables}
\else
  \newcommand\listtablename{List of Tables}
\fi
\ifdefined\figurename
  \renewcommand*\figurename{Figure}
\else
  \newcommand\figurename{Figure}
\fi
\ifdefined\tablename
  \renewcommand*\tablename{Table}
\else
  \newcommand\tablename{Table}
\fi
}
\@ifpackageloaded{float}{}{\usepackage{float}}
\floatstyle{ruled}
\@ifundefined{c@chapter}{\newfloat{codelisting}{h}{lop}}{\newfloat{codelisting}{h}{lop}[chapter]}
\floatname{codelisting}{Listing}
\newcommand*\listoflistings{\listof{codelisting}{List of Listings}}
\usepackage{amsthm}
\theoremstyle{definition}
\newtheorem{definition}{Definition}[section]
\theoremstyle{plain}
\newtheorem{theorem}{Theorem}[section]
\theoremstyle{remark}
\AtBeginDocument{\renewcommand*{\proofname}{Proof}}
\newtheorem*{remark}{Remark}
\newtheorem*{solution}{Solution}
\newtheorem{refremark}{Remark}[section]
\newtheorem{refsolution}{Solution}[section]
\makeatother
\makeatletter
\makeatother
\makeatletter
\@ifpackageloaded{caption}{}{\usepackage{caption}}
\@ifpackageloaded{subcaption}{}{\usepackage{subcaption}}
\makeatother
\usepackage{bookmark}
\IfFileExists{xurl.sty}{\usepackage{xurl}}{} % add URL line breaks if available
\urlstyle{same}
\hypersetup{
  pdftitle={Test Article for Custom AmSThm Environments},
  pdfauthor={Test Author},
  colorlinks=true,
  linkcolor={blue},
  filecolor={Maroon},
  citecolor={Blue},
  urlcolor={Blue},
  pdfcreator={LaTeX via pandoc}}


\title{Test Article for Custom AmSThm Environments}
\author{Test Author}
\date{2025-10-14}
\begin{document}
\maketitle


\section{Introduction}\label{introduction}

This is a minimal test article to demonstrate the
custom-amsthm-environments extension in a single-document format.

\begin{problem}[First Problem]\label{prm-first}

Verify that this extension works correctly.

\end{problem}

As stated in Problem~\ref{prm-first}, we need to ensure proper
functionality.

\section{Section with Multiple
Environments}\label{section-with-multiple-environments}

Here we test different types of custom environments:

\begin{problem}\label{prm-second}

This problem doesn't have a title.

\end{problem}

\begin{task}[Key Task]\label{tsk-key}

For all \(x \in \mathbb{R}\), we have \(x^2 \geq 0\).

\end{task}

\begin{use case}

This is just a unnumbered use case description

\end{use case}

References: Problem~\ref{prm-second}, Task~\ref{tsk-key} work in the
text.

\section{Testing Numbering}\label{testing-numbering}

With section-based numbering, problems in different sections should have
different prefixes:

\begin{problem}\label{prm-subsection1}

Problem in section 2.2.

\end{problem}

See Problem~\ref{prm-subsection1}.

\section{Subsection Level}\label{subsection-level}

\begin{problem}[Deeper Problem]\label{prm-subsection2}

This problem is in section 2.2.1.

\end{problem}

Global numbering (tasks) should just increment:

\begin{task}\label{tsk-second}

Another task with global numbering.

\end{task}

Compare Task~\ref{tsk-key} (Task 1) with Task~\ref{tsk-second} (Task 2).

\section{Cross-References Summary}\label{cross-references-summary}

Let's verify all our cross-references work:

\begin{enumerate}
\def\labelenumi{\arabic{enumi}.}
\tightlist
\item
  Problems (section-based): Problem~\ref{prm-first},
  Problem~\ref{prm-second}, Problem~\ref{prm-subsection1},
  Problem~\ref{prm-subsection2}
\item
  Tasks (global): Task~\ref{tsk-key}, Task~\ref{tsk-second}
\item
  Unnumbered use cases cannot be referenced reliably
  (Use Case~\ref{usecase-observation})
\end{enumerate}

\section{Built-in Environments}\label{built-in-environments}

The extension should coexist with Quarto's standard environments:

\begin{theorem}[Standard
Theorem]\protect\hypertarget{thm-standard}{}\label{thm-standard}

This uses Quarto's built-in theorem support.

\end{theorem}

\begin{definition}[]\protect\hypertarget{def-standard}{}\label{def-standard}

A \textbf{widget} is a conceptual object.

\end{definition}

Both Theorem~\ref{thm-standard} and Definition~\ref{def-standard} should
work alongside our custom environments like Problem~\ref{prm-first} and
Task~\ref{tsk-key}.




\end{document}
